\documentclass[10pt,a4paper]{amsart}
\usepackage[margin=19mm]{geometry}
\usepackage{amsmath,amssymb,mathtools}
\usepackage[scr=rsfs,cal=euler]{mathalfa}
\usepackage{xfrac}
\usepackage[unicode,hidelinks,bookmarks=false]{hyperref}
\newcommand{\ie}{i.e.\ }
\newcommand{\cf}{cf.\ }
\newcommand{\resp}{respectively\ }
\newcommand{\Subsection}{\subsection}
\providecommand{\nobreakdash}{\nobreak\hbox{-}}
\setlength{\emergencystretch}{2em}
\multlinegap0pt
\usepackage[shortlabels]{enumitem}
\usepackage{amssymb}
\usepackage{mathtools}
\usepackage{tikz}
\newtheorem{lemma}{Lemma}
\DeclareMathOperator{\ccl}{ccl}
\newcommand{\eps}{\varepsilon}
\title{\bf Large Complete Minors from a Cheeger Condition}
\author{Chengli Li, Leyou Xu, Bo Zhou}
\date{Source: arXiv:2609.15414v1 (CC BY 4.0)}
\begin{document}
\maketitle

\noindent\emph{Source and licence.} \bf Large Complete Minors from a Cheeger Condition. Version arXiv:2609.15414v1 (CC BY 4.0), distributed by arXiv under the Creative Commons Attribution 4.0 International licence, http://creativecommons.org/licenses/by/4.0/, as recorded in the arXiv licence metadata for this exact version and shown on the abstract page https://arxiv.org/abs/2609.15414v1. Full licence notice: \texttt{licenses/arxiv-2609.15414-CC-BY-4.0.txt}.\par\smallskip

\noindent\textit{Editorial scope.} Counted unit: the article's lemma lem:scale (source0.tex lines 993--1005). The article's complete original proof of it is delivered with it (source0.tex lines 1007--1312, one \texttt{proof} environment of 306 lines). No mathematics is written, repaired or re-derived: the statement and the proof are the article's own lines, in the article's own notation, and no retained line is substituted or re-flowed.  Retained uncounted as the source-local context the proof needs: lines 380--410; lines 553--579; lines 884--904.  Nothing was omitted from any retained span, so the package carries no omission record.  No retained line contains a citation, so the wrapper carries no bibliography and no bibliography entry.  No retained line contains a figure environment, an image-inclusion command or drawing code, so no image is delivered and the package declares no asset dependency.  Editorial formatting only, in addition: an AMS \texttt{amsart} preamble that restates the article's own declarations and macros used by the retained lines, namely the environments \texttt{lemma} on separate counters, together with the standard packages those lines need.  The retained environments print in this document as Lemma 1 in order of appearance, whereas the article prints the same items with the numbering of its own full document.  The complete native source of the article as distributed in the arXiv e-print for this exact version is bundled unchanged as \texttt{sources/210414/source0.tex}.
\begin{lemma}\label{lem:remainder}
Let $0<\eps\le1$, let $d>0$, and let $G$ be a graph on $n$ vertices
satisfying $\Delta(G)\le d$ and $h(G)\ge\eps d$. Suppose that
\[
 V(G)=Y\mathbin{\dot\cup}D_0\mathbin{\dot\cup}U_0,
 \ |Y|\le\frac{\eps^2n}{16},
 \ |D_0|\le\frac{n}{2}
 \  \text{ and }\ e_G(D_0,U_0)\le\frac{\eps d}{4}|D_0|.
\]
Then there is a partition
$V(G)=Y\mathbin{\dot\cup}D\mathbin{\dot\cup}U$ such that
$D_0\subseteq D$ and
\[
 e_G(D,U)\le\frac{\eps d}{4}|D|,
 \ |D|\le\frac{4|Y|}{3\eps},
 \ |U|\ge\frac{n}{2}
 \ \text{ and } \  h(G[U])\ge\frac{\eps d}{4}.
\]
Moreover, if $Z\subseteq U$ and
$|Y\cup Z|\le\frac{\eps^2n}{16}$, then there is a partition
$
 V(G)=(Y\cup Z)\mathbin{\dot\cup}D'\mathbin{\dot\cup}U'
$
such that $D\subseteq D'$, $U'\subseteq U\setminus Z$, and
\[
 e_G(D',U')\le\frac{\eps d}{4}|D'|,
 \ |D'|\le\frac{4|Y\cup Z|}{3\eps},
 \ |U'|\ge\frac{n}{2}
 \ \text{ and }\  h(G[U'])\ge\frac{\eps d}{4}.
\]
\end{lemma}

\begin{lemma}\label{lem:stable-branches}
For every $\eps\in(0,1]$, there are constants
$c=c(\eps)>0$, $\xi=\xi(\eps)>0$, $\zeta=\zeta(\eps)>0$, and an
integer $d_2=d_2(\eps)$ such that 
\[
c\le\frac{1}{32}, \
\xi\le\frac{1}{2048}\ \text{ and } \ \zeta\le\eps,
\]
and
the following holds. Let $G$ be a
graph on $n$ vertices satisfying $\Delta(G)\le d$,
$h(G)\ge\eps d$, and $d\ge d_2$. If $\ell$ is an integer satisfying
$2\le\ell\le\frac{n}{2d}$, then there exist an integer $m\ge1$, pairwise
disjoint nonempty connected sets $A_1,\ldots,A_m$, and sets
$Q_1,\ldots,Q_m$ with the following properties. If
$X=\bigcup_{i=1}^mA_i$, then
\begin{enumerate}[label=\textup{(\roman*)}]
 \item $\frac{\xi n}{\ell}\le m\le\frac{4\xi n}{\ell}$ and
       $|X|\le\frac{\eps\zeta n}{32}$;
 \item $|A_i|\le\ell$, $Q_i\subseteq N_G(A_i)\setminus X$, and
       $|Q_i|\ge cd\ell$ for every $i\in[m]$;
 \item for every $Z\subseteq V(G)$ with $|Z|\le\zeta n$, fewer than
       $\frac{m}{16}$ indices $i\in[m]$ satisfy
       $|Q_i\setminus Z|<cd\ell$.
\end{enumerate}

\end{lemma}

\begin{lemma}\label{lem:scale-hitting}
For every $\alpha\in(0,\frac{1}{2})$, $c\in(0, \frac{1}{32}]$ and
$\xi\in(0,\frac{1}{2048}]$, there is a constant $C=C(\alpha,c)>0$ and an
integer $d_3=d_3(\alpha,c,\xi)\ge3$ such that the following holds.
Let $n,d,\ell$ be positive integers satisfying $d\ge d_3$ and
$2\le\ell\le \frac{n}{2d}$, and let $H$ be a graph on $N$ vertices with
$\frac{n}{2}\le N\le n$, $\Delta(H)\le d$ and $h(H)\ge\alpha d$.
Suppose that $q$ is a positive integer satisfying
$\frac{15\xi n}{16\ell}\le q\le \frac{4\xi n}{\ell}$, and that
$W_1,\ldots,W_q\subseteq V(H)$ satisfy $|W_i|\ge cd\ell$ for
every $i\in[q]$.
\begin{enumerate}[label=\textup{(\alph*)}]
 \item There is a connected set $T\subseteq V(H)$ which
       meets every $W_i$ and satisfies
       $|T|\le \frac{Cn\log d}{d\ell}$.
 \item There is a connected set $T\subseteq V(H)$ which
       meets at least $\frac{3}{4}q$ of the sets $W_1,\ldots,W_q$ and
       satisfies $|T|\le \frac{Cn}{d\ell}$.
\end{enumerate}
%The sets in \textup{(a)} and \textup{(b)} may be different.
\end{lemma}

\begin{lemma}\label{lem:scale}
For every $\eps\in(0,1]$, there exist a constant $b=b(\eps)>0$ and an
integer $d_0=d_0(\eps)$ such that the following holds. Let $G$ be a
graph on $n$ vertices, let $d\ge d_0$ be an integer, and suppose that
$\Delta(G)\le d$ and $h(G)\ge\eps d$. If $\ell$ is an integer
satisfying $2\le\ell\le\frac{n}{2d}$, then the following statements hold.
\begin{enumerate}[label=\textup{(\roman*)}]
 \item If $d\ell^2\ge n\log d$, then
       $\ccl(G)\ge\frac{bn}{\ell}$.
 \item If $d\ell^2\ge n$, then $G$ has a minor of average degree at
       least $\frac{bn}{\ell}$.
\end{enumerate}
\end{lemma}

\begin{proof}
Let $c,\xi,\zeta$ and $d_2$ be given in
Lemma~\ref{lem:stable-branches}.  In particular, $c\le\frac{1}{32}$,
$\xi\le\frac{1}{2048}$ and $\zeta\le\eps$. Let $\alpha=\eps/4$.
Then let $C,d_3$ be given in  Lemma~\ref{lem:scale-hitting}.


Choose $\eta>0$ such that
\[
 \eta\le\frac{1}{16}
 \ \text{ and }\ 
 4\eta\xi C\le\frac{\eps\zeta}{32}.
\]
Choose  an integer $d_0\ge \max\{d_2, d_3\}$ such that  $\eta\xi d_0\ge1$. 
All these choices depend only on $\eps$. Note that $d\ge d_0$. 


By Lemma~\ref{lem:stable-branches}, for some integer $m\ge 1$, there are 
pairwise disjoint nonempty connected sets $A_1,\ldots,A_m$ and
sets $Q_1,\ldots,Q_m$ such that 
\[
 \frac{\xi n}{\ell}\le m\le\frac{4\xi n}{\ell},
 \ |X|\le\frac{\eps\zeta n}{32}
 \ \text{ and } \  Q_i\subseteq N_G(A_i)\setminus X
 \text{ for } i\in[m]
\]
with 
\[
 X=\bigcup_{i=1}^mA_i.
\]
Let $k=\lfloor\eta m\rfloor$. Since $m\ge\xi n/\ell$,
$n/\ell\ge2d$ and $\eta\xi d_0\ge1$, we have
$\eta m\ge2\eta\xi d\ge2$, so
\[
1\le\frac{\eta m}{2}\le k\le\eta m\le\frac m{16}.
\]

For a set $U\subseteq V(G)\setminus X$, call an index $i\in[m]$
\emph{$U$-active} if $|Q_i\cap U|\ge cd\ell$.




\medskip
\noindent\textbf{Claim.}
For each of parts (i) and (ii), under the
corresponding hypothesis, there exist pairwise disjoint connected sets
$T_j$, $j\in[k]$, all disjoint from $X$, and sets
$D_0,\ldots,D_k,U_0,\ldots,U_k$ (these sets
may differ between part (i) and part (ii)) with the following properties. 
For $0\le j\le k$, let
\[
 F_0=\emptyset
 \text{ and }
 F_j=T_1\mathbin{\dot\cup}\cdots\mathbin{\dot\cup}T_j
 \text{ for } j\ge 1.
\]
For every $j$ with $0\le j\le k$,
\begin{equation}\label{eq:scale-recursive-partition}
 V(G)=(X\cup F_j)\mathbin{\dot\cup}D_j
       \mathbin{\dot\cup}U_j,
\end{equation}
\[
 e_G(D_j,U_j)\le\alpha d|D_j|,
 \ 
 |D_j|\le\frac{4|X\cup F_j|}{3\eps},
 \ 
 |U_j|\ge\frac{n}{2}, \
 h(G[U_j])\ge\alpha d,
\]
and at least $\frac{15m}{16}$ indices are $U_j$-active.
For every $j\in [k]$,
\[
 T_j\subseteq U_{j-1},
 \ 
 D_{j-1}\subseteq D_j,
 \ 
 U_j\subseteq U_{j-1}\setminus T_j \
 \text{ and }\ 
 |T_j|\le C\ell.
\]
$T_j$ meets $Q_i\cap U_{j-1}$
for every $U_{j-1}$-active index $i$ in part (i), 
and 
$T_j$ meets $Q_i\cap U_{j-1}$ for at least $\frac34$ of
the $U_{j-1}$-active indices $i$ in part (ii).

\begin{proof}
We construct the desired sets recursively.
For the initial stage, let $F_0=\emptyset$. Recall that
$|X|\le\frac{\eps\zeta n}{32}\le\frac{\eps^2n}{16}$.
So we can apply 
Lemma~\ref{lem:remainder} to the partition
\[
 V(G)=X\mathbin{\dot\cup}\emptyset
       \mathbin{\dot\cup}(V(G)\setminus X)
\]
to obtain the new partition $V(G)=X\mathbin{\dot\cup}D_0
       \mathbin{\dot\cup}U_0$ such that
\[
 e_G(D_0,U_0)\le\alpha d|D_0|,
 \ 
 |D_0|\le\frac{4|X\cup F_0|}{3\eps}, \ 
 |U_0|\ge\frac{n}{2} \ 
 \text{ and } \
 h(G[U_0])\ge\alpha d.
\]
As $|X|\le\frac{\eps\zeta n}{32}$, we have
\[
 |D_0|\le\frac{4|X|}{3\eps}
 \le\frac{\zeta n}{24}<\zeta n.
\]
By Lemma \ref{lem:stable-branches} (ii), $Q_i\cap X=\emptyset$ for every $i\in[m]$. So, from $V(G)=X\mathbin{\dot\cup}D_0
       \mathbin{\dot\cup}U_0$, we have
\[
 Q_i\setminus D_0=Q_i\cap U_0
\ \text{ for }i\in[m].
\]
Now Lemma~\ref{lem:stable-branches} (iii) applied
with $Z=D_0$  shows that at least $\frac{15m}{16}$ indices are
$U_0$-active. Thus the claim holds for $j=0$.


Now suppose that the claim holds  for $0\le j<k$. Let $q_j$ be the number of $U_j$-active indices. By induction hypothesis and the upper bound on $m$, we have
\[
 \frac{15\xi n}{16\ell}
 \le\frac{15m}{16}
 \le q_j\le m\le\frac{4\xi n}{\ell}.
\]
Furthermore, $\frac{n}{2}\le |U_j|\le n$, $\Delta(G[U_j])\le d$ and
$h(G[U_j])\ge\alpha d$. Let $W_1,\ldots,W_{q_j}$ be the sets $Q_i\cap U_j$
for $i\in[m]$ such that $i$ is $U_j$-active. Then $|W_i|\ge cd\ell$ for each $i\in [q_j]$. 
Note that $\alpha\in(0,1/2)$, $d\ge d_0\ge d_3$ and $2\le\ell\le n/(2d)$. 
So all the hypotheses of Lemma~\ref{lem:scale-hitting} for
$H=G[U_j]$ are satisfied with $q=q_j$.





In part (i), Lemma \ref{lem:scale-hitting} (a) ensures that there is a connected set $T_{j+1}\subseteq U_j$ meeting $Q_i\cap U_j$ for every $U_j$-active index $i$. Since $d\ell^2\ge n\log d$, we have 
$|T_{j+1}|\le \frac {Cn\log d}{d\ell}\le C\ell$; 
In part (ii), Lemma \ref{lem:scale-hitting} (b) ensures that there is a connected set $T_{j+1}$ meeting $Q_i\cap U_j$ for at least $\frac{3}{4}$ of the
$U_j$-active indices. Since $d\ell^2\ge n$, we have 
$|T_{j+1}|\le\frac{Cn}{d\ell}\le C\ell$.
Thus $T_{j+1}$ has the required properties in either part.


Since $T_{j+1}\subseteq U_j$, we have from 
\eqref{eq:scale-recursive-partition} that  the set
$T_{j+1}$ is disjoint from $X\cup F_j\cup D_j$, so 
$T_1,\ldots,T_{j+1}$ are pairwise disjoint and are disjoint from
$X$. Define
\[
 F_{j+1}=F_j\mathbin{\dot\cup}T_{j+1}.
\]
Since every set $T_r$, $1\le r\le j+1$, was supplied by Lemma \ref{lem:scale-hitting}, we have $|T_r|\le C\ell$,
so
\begin{equation}\label{eq:scale-F-bound}
 \begin{aligned}
 |F_{j+1}|
 \le (j+1)C\ell
 \le kC\ell
 \le\eta mC\ell
 \le4\eta\xi Cn
 \le\frac{\eps\zeta n}{32},
 \end{aligned}
\end{equation}
so
\begin{equation}\label{MM}
 |X\cup F_{j+1}|\le |X|+|F_{j+1}|
 \le\frac{\eps\zeta n}{16}
 \le\frac{\eps^2n}{16}.
\end{equation}
Recall that $\zeta\le\eps\le1$. From the induction hypothesis and the fact that 
$X\cup F_j\subseteq X\cup F_{j+1}$, we have
\[
 |D_j|
 \le\frac{4|X\cup F_j|}{3\eps}
 \le\frac{4|X\cup F_{j+1}|}{3\eps}
 \le\frac{\zeta n}{12}<\frac{n}{2}.
\]
Moreover,
\[
 e_G(D_j,U_j\setminus T_{j+1})
 \le e_G(D_j,U_j)
 \le\alpha d|D_j|.
\]
Thus Lemma~\ref{lem:remainder} applies to the partition
\[
 V(G)=(X\cup F_{j+1})\mathbin{\dot\cup}D_j
       \mathbin{\dot\cup}(U_j\setminus T_{j+1})
\]
to form the new partition 
\[
 V(G)=(X\cup F_{j+1})\mathbin{\dot\cup}D_{j+1}
       \mathbin{\dot\cup}U_{j+1}
\]
with 
$ D_j\subseteq D_{j+1}$,  
$U_{j+1}\subseteq U_j\setminus T_{j+1}$,
\[
 e_G(D_{j+1},U_{j+1})\le\alpha d|D_{j+1}|,
 \ 
 |D_{j+1}|\le\frac{4|X\cup F_{j+1}|}{3\eps},
\  
 |U_{j+1}|\ge\frac{n}{2}
 \  \text{ and }\ 
 h(G[U_{j+1}])\ge\alpha d.
\]
%It remains only to verify that sufficiently many indices are
%$U_{j+1}$-active. 
From $|D_{j+1}|\le\frac{4|X\cup F_{j+1}|}{3\eps}$ and \eqref{MM}, we have 
\[
 |D_{j+1}|\le\frac{\zeta n}{12},
\]
which, together with
\eqref{eq:scale-F-bound} gives
 \[
 |D_{j+1}\cup F_{j+1}|
 \le\frac{\zeta n}{12}+\frac{\eps\zeta n}{32}
 <\zeta n.
\]
Since $Q_i\cap X=\emptyset$ and $V(G)=(X\cup F_{j+1})\mathbin{\dot\cup}D_{j+1}
       \mathbin{\dot\cup}U_{j+1}$,
\[
 Q_i\setminus(D_{j+1}\cup F_{j+1})=Q_i\cap U_{j+1}
\ \text{ for every }i\in[m].
\]
Now property (iii) of Lemma~\ref{lem:stable-branches} applied
with $Z=D_{j+1}\cup F_{j+1}$ shows that at least
$\frac{15m}{16}$ indices are $U_{j+1}$-active. This completes the
inductive step and hence the construction through stage $k$.
\end{proof} 
%This proves Claim~2. %\hfill$\square$

Set $b=\frac{45\eta\xi}{128}$, which depends only on $\eps$.
We now prove parts (i) and (ii) with this choice of $b$.

%\medskip
%\noindent\emph{Proof of part (i).}
Suppose first that $d\ell^2\ge n\log d$.  We use the sets constructed
in the  Claim for part (i). At least $\frac{15m}{16}$ indices are
$U_k$-active.  Since  $k\le\frac{m}{16}$, we can choose distinct
$U_k$-active indices $i_1,\ldots,i_k$.

Fix $j,r\in[k]$. Since $U_k\subseteq U_{j-1}$, we have
\[
 |Q_{i_r}\cap U_{j-1}|
 \ge |Q_{i_r}\cap U_k|
 \ge cd\ell.
\]
Thus $i_r$ is $U_{j-1}$-active. By the 
Claim (for part (i)), $T_j$ meets $Q_{i_r}\cap U_{j-1}$.
As $Q_{i_r}\subseteq N_G(A_{i_r})$, there is an edge between
$T_j$ and $A_{i_r}$.

For each $j\in[k]$, let $B_j=A_{i_j}\cup T_j$.
The sets $B_1,\ldots,B_k$ are pairwise disjoint, because the
sets $A_i$ are pairwise disjoint, the sets $T_j$ are pairwise
disjoint, and every $T_j$ is disjoint from $X$.
Since $A_{i_j}$ and $T_j$ are connected
and there is an edge between them, we deduce that each $B_j$ is connected. Moreover, for distinct
$j,r\in[k]$, the edge between $T_j$ and $A_{i_r}$ joins $B_j$
to $B_r$. Hence $(B_j)_{j\in[k]}$ is a $K_k$-minor model in $G$.
Using $m\ge\frac{\xi n}{\ell}$,
we obtain
\[
 \ccl(G)\ge k
 \ge\frac{\eta m}{2}
 \ge\frac{\eta\xi n}{2\ell}
 \ge\frac{bn}{\ell},
\]
where the last inequality follows from $\frac{45}{128}<\frac12$.
This proves part (i).

%\medskip
%\noindent\emph{Proof of part (ii).}
Suppose next that $d\ell^2\ge n$, and use the sets constructed in the
Claim for part (ii). Define a bipartite graph $J$ with parts
$\{A_1,\ldots,A_m\}$ and $\{T_1,\ldots,T_k\}$, joining $A_i$
to $T_j$ precisely when $e_G(A_i,T_j)>0$.
All these vertex sets are pairwise disjoint and connected, so
they form a $J$-minor model in $G$. Thus $J\prec G$.

For each $j\in[k]$, at least $\frac{15m}{16}$ indices are
$U_{j-1}$-active. The meeting property in the Claim for part (ii)
therefore ensures that $T_j$ meets $Q_i\cap U_{j-1}$ for at least
$\frac34\cdot\frac{15m}{16}=\frac{45m}{64}$ distinct indices $i$.
Each such intersection gives an edge between $T_j$ and $A_i$,
since $Q_i\subseteq N_G(A_i)$. Consequently,
\[
 |E(J)|=\sum_{j=1}^k\deg_J(T_j)
 \ge\frac{45km}{64}.
\]
Since $|V(J)|=m+k\le2m$, $k\ge \eta \frac{m}{2}$ and $m\ge \frac{\xi n}{\ell}$, we have
\[
 \overline d(J)
 =\frac{2|E(J)|}{m+k}
 \ge\frac{45k}{64}
 \ge\frac{45\eta m}{128}
 \ge\frac{45\eta\xi n}{128\ell}
 =\frac{bn}{\ell}.
\]
This proves part (ii).
\end{proof}

\end{document}
